%% kapa: exam=1 round=32 year=2021 period=1 form=A subject=civil_law subject_ko=민법 seq=13 number=13
%% source: Q-Net 감정평가사 기출문제 (cst00309) artlSeq=5208397 file=extracted/1008/5208397/A형/2021년도 제32회 감정평가사 1차 1교시 A형.pdf
%% conversion: kordoc 4.24.51 -> tools/gen1.py (auto); review=proofread-sonnet; flags=none
\begin{question}{13}
사기·강박에 의한 의사표시에 관한 설명으로 옳은 것은? (다툼이 있으면 판례에 따름)
\begin{choices}
  \item 교환계약의 당사자가 자기 소유 목적물의 시가를 묵비하였다면 특별한 사정이 없는 한, 위법한 기망행위가 성립한다.
  \item 강박에 의해 자유로운 의사결정의 여지가 완전히 박탈되어 그 외형만 있는 법률행위라고 하더라도 이를 무효라고 할 수는 없다.
  \item 토지거래허가를 받지 않아 유동적 무효 상태에 있는 법률행위라도 사기에 의한 의사표시의 요건이 충족된 경우 사기를 이유로 취소할 수 있다.
  \item 대리인의 기망행위로 계약을 체결한 상대방은 본인이 대리인의 기망행위에 대해 선의·무과실이면 계약을 취소할 수 없다.
  \item 강박행위의 목적이 정당한 경우에는 비록 그 수단이 부당하다고 하더라도 위법성이 인정될 여지가 없다.
\end{choices}
\end{question}
\answer{3}
